% ------------------------------------------------------------- Appendix
\section{Notation and Acronyms}
\label{app:notation}
\begin{center}\footnotesize
\begin{tabularx}{\columnwidth}{@{}l>{\raggedright\arraybackslash}X@{}}
\toprule
\textbf{Symbol} & \textbf{Meaning} \\
\midrule
$x,\hat x$ & true and estimated navigation state \\
$e(t)$ & cross-track deviation from the nominal path (m) \\
$m$ & corridor margin (m); completion iff $\sup_t|e|\le m$ \\
$L,\ T_c$ & closed-loop Lipschitz constant; certification horizon (1\,s) \\
$\theta,\ \theta_D$ & detector operating point; threshold of detector $D$ \\
$\Delta(\theta),\ s,\ H$ & hidden staleness; contact-window slack; malicious hops \\
$\gamma$ & delay-to-position rate (m/s), campaign supremum \\
$\dnet,\dadv$ & deterministic network and adversarial-offset budgets (m) \\
$\dgnss,\dfault,\denv$ & stochastic and bounded-unknown budget terms \\
$X_i$ & sup-deviation of the closed loop under source $i$ alone \\
$S_i$ & survival function of $X_i$, $S_i(t)=\Pr[X_i>t]$ \\
$\underline F,\overline F$ & lower and upper CDFs of a p-box ($\underline F$ is stochastically largest) \\
$\pbar,\ p_U$ & worst-case (all couplings) and union (standard) tail bounds \\
$f(m)$ & certified floor, $1-\pbar(s_m)$ \\
$U$ & event that the detector stack raises no alarm \\
$T_a$ & re-anchoring interval (C2$'$) \\
$O_{0:T}$ & defender's observations over the mission \\
$P_a,\ P_d$ & laws of $O_{0:T}$ under attack $a$ / disguise $d$ \\
$\Dclass_k$ & disguise family $k$ \\
\bottomrule
\end{tabularx}

\vspace{4pt}
\begin{tabularx}{\columnwidth}{@{}l>{\raggedright\arraybackslash}X@{}}
\toprule
\textbf{Acronym} & \textbf{Expansion} \\
\midrule
CDF & cumulative distribution function \\
CEM & cross-entropy method \\
CMA-ES & covariance matrix adaptation evolution strategy \\
CUSUM & cumulative sum (sequential change detection) \\
EKF / KF & extended Kalman filter / Kalman filter \\
FPR & false-positive rate \\
GNSS & global navigation satellite system \\
IDS & intrusion detection system \\
IMU & inertial measurement unit \\
KL / KS & Kullback--Leibler divergence / Kolmogorov--Smirnov distance \\
LoE & loss of effectiveness (actuator) \\
MCR & mission-completion rate \\
NIS & normalised innovation squared \\
PD & proportional--derivative (controller) \\
p-box & probability box \\
RA & Rearrangement Algorithm \\
SITL & software in the loop \\
TV & total-variation distance \\
UAV & unmanned aerial vehicle \\
\bottomrule
\end{tabularx}
\end{center}

\section{Proof Details}
\label{app:proofs}
\paragraph{Proof of \cref{prop:union}.}
(a) For any split, $\{\sum X_i>s\}\subseteq\bigcup_i\{X_i>s_i\}$. (b) is
Makarov's theorem~\cite{makarov1981}, proved independently by
R\"uschendorf~\cite{ruschendorf1982}; Frank et al.\ give the
attaining coupling~\cite{frank1987}. (c) Embrechts et al.\ give examples of
strict inequality~\cite{embrechts2013var}. (d) is direct
computation; for independent exponentials, $p_U$ is attained at the equal split
and the ratio $2e^{s/2}/(1+s)$ diverges (\cref{app:proofs}). (e) The Fr\'echet
class of a Dirac marginal is a singleton.

\paragraph{Proof of \cref{thm:floor}.}
On $U$, (C3) gives $\sup_t|e_{\mathrm{adv}}|\le\dadv$, and \eqref{eq:dnet} bounds
the network contribution by $e^{LT_c}\dnet(\theta)$, which adds to the others
by (C6) (\cref{lem:subadd} does not cover the network term). By \cref{lem:subadd} and
(C5), failure on $U$ implies $\sum_iX_i>s_m$. Hence
$\Pr[\text{fail}\wedge U]\le\Pr[\sum_iX_i>s_m]\le\pbar(s_m)$, which needs no
conditioning on $U$. The tail of a sum is non-decreasing under first-order
stochastic dominance of any summand, for every coupling and hence for the
supremum, so the stochastically largest member of the p-box is the worst
(\cref{app:proofs}). The conditional form follows by dividing by $\Pr[U]$.

\paragraph{Proof of \cref{prop:certbudget}.}
By superposition (\cref{lem:subadd}, exact on the event $U$ because
$D_{\mathrm{sat}}$ is silent), the attack-only deviation and the attack-only parts
of the detector statistics are linear in $d$. Silence of each detector in $\Psi$, together with
the benign margins, implies a convex (here polyhedral) constraint on $d$. The
bound holds for every $d$ in the feasible set, which contains every sequence an
adaptive policy can produce on that realisation. The order statistic
$X_{(k)}$ falls below the $q$-quantile only if at least $k$ of the $N$ samples do,
with probability $P(\mathrm{Bin}(N,q)\ge k)\le\delta/m$ for any continuous
distribution; union bounds over the $m$ margins give the confidence and the
probability.

\paragraph{\Cref{prop:union}(b): sharpness for two sources.}
For $d=2$, Makarov~\cite{makarov1981} showed that
$\sup_C\Pr_C[X_1+X_2>s]=\min\{1,\inf_{t}(S_1(t^-)+S_2((s-t)^-))\}$,
which equals $p_U(s)$ at every continuity point of the empirical survival
functions. The extremal coupling makes the events $\{X_1>t^\star\}$ and
$\{X_2>s-t^\star\}$ disjoint, by pairing the upper tail of one source with the
body of the other~\cite{frank1987}. An attacker who controls when a stochastic
source fires (for example, the timing of jamming-induced outages relative to
gusts) can move the coupling towards this extremal one.

\paragraph{\Cref{prop:union}(d): comonotone and independent gaps.}
Let $X_1=\dots=X_d=X$ almost surely, with survival function $S$. Then
$\Pr[\sum X_i>s]=S(s/d)$, while the equal split gives $d\,S(s/d)$; whenever the
equal split is optimal, the ratio is $d$. For two independent $\mathrm{Exp}(1)$
variables, $\Pr[X_1+X_2>s]=(1+s)e^{-s}$, while
$\inf_t\,(e^{-t}+e^{-(s-t)})=2e^{-s/2}$. The ratio $2e^{s/2}/(1+s)$ is unbounded.

\paragraph{Monotonicity in the p-box (proof of \cref{thm:floor}).}
If $X_e'\succeq_{\mathrm{st}}X_e$, there is a coupling with $X_e'\ge X_e$ almost
surely. For any copula of $(X_g,X_f,X_e)$, compose it with this monotone
coupling to obtain a copula of $(X_g,X_f,X_e')$ under which the sum is
pointwise no smaller. Hence $\pbar$ evaluated at the stochastically largest
member of the box (CDF $\underline F_{\mathrm{env}}$; in the pilot, the 30\,kt
marginal) dominates $\pbar$ at every $F_{\mathrm{env}}$ in the box.

\section{Rearrangement Algorithm}
\label{app:ra}
For a confidence level $\alpha$, build an $N\times d$ matrix whose column $j$
holds the quantiles of $X_j$ at the levels $\alpha+(1-\alpha)\,i/N$, for $i=1,\dots,N$
in the upper discretisation and $i=0,\dots,N-1$ in the lower one. Iterate until
the minimum row sum stabilises. In each iteration, rearrange every column so
that it is oppositely ordered to the sum of the other columns. The resulting
minimum row sum approximates the worst-case value-at-risk
$\overline{\mathrm{VaR}}_\alpha(\sum X_j)$~\cite{embrechts2013var}, and the
conservative of the two discretisations is kept. Then
$\pbar(s)\le1-\sup\{\alpha:\overline{\mathrm{VaR}}_\alpha\le s\}$, found by
bisection. The RA's convergence to the sharp bound is not proven in general.
For $d=2$, our implementation (\code{depbounds.py}) reproduces $\pbar=p_U$ to
within \nUnionRaDtwoMaxErr{} wherever $p_U>0.002$; below that level, the RA's
$1/N$ resolution floor dominates, which is why we report the minimum of the two
bounds.

\section{Simulator Constants}
\label{app:constants}
\begin{center}\footnotesize
\begin{tabularx}{\columnwidth}{@{}l l >{\raggedright\arraybackslash}X@{}}
\toprule
\textbf{Constant} & \textbf{Value} & \textbf{Provenance} \\
\midrule
$\gamma$ & \nParentGamma\,m/s & measured, UAV Attack Dataset~\cite{whelan2020dataset} (674 samples) \\
$s$ & \nParentSlack\,s & measured, artifact (\nParentHops{} hops) \\
$L$ & \nParentL & artifact, local estimate \\
step / GNSS rate & 20 / 10\,Hz & model \\
PD gains $(k_p,k_d)$ & (1.2, 1.8) & model \\
drag $c_d$ & 0.35\,s$^{-1}$ & model \\
Dryden $h$, $V$ & 100\,ft, 8\,m/s & MIL-F-8785C \\
$W_{20}$ envelope & 15--30\,kt & declared (C4) \\
GNSS white / GM & 0.5\,m / 0.8\,m, 60\,s & model \\
accel.\ noise / bias RW & 0.05\,m/s$^2$ / 0.002\,m/s$^2/\sqrt{\mathrm s}$ & model (Allan class~\cite{elsheimy2008}) \\
anchor $\sigma$, $T_a$ & 0.3\,m, 5\,s & model (C2$'$) \\
fault probabilities & 0.3 / 0.2 / 0.2 & model (RflyMAD categories) \\
outage / jam & 0.3 (mean 5\,s) / 0.2 & model \\
\bottomrule
\end{tabularx}
\end{center}

\section{Looseness of the Bound (RQ3)}
\label{app:rq3}
\Cref{fig:gap} decomposes the tail of the benign budget; the ratios below are
at $s=\nGapS$\,m, and their ordering changes in the far tail, where few joint
exceedances remain.
\begin{enumerate}
\item For two sources, the union and RA bounds agree to within
  \nUnionRaDtwoMaxErr{}, as Makarov's theorem requires. For three, the RA
  tightens the union bound by \nRaGain{} (\nGapUnion{} to \nGapRa{}). Under
  unknown dependence, the union bound is nearly the right quantity.
\item The worst coupling exceeds independent sums at the same marginals by
  \nDepRatio$\times$ (\emph{dependence gap}). Evaluating the wind at the upper
  envelope rather than at the conditions actually flown costs a further
  \nEnvRatio$\times$ (\emph{envelope gap}).
\item Independent sums at matched wind agree with the sum of per-source
  suprema measured on the same missions to within \nNatRatio$\times$,
  consistent with the benign sources being independent, as they are by
  construction in the simulator.
\item The sum of per-source suprema exceeds the supremum of the joint deviation
  by \nTimeRatio$\times$ (\emph{temporal gap}; \nGapJointK{} joint exceedances in
  \nJointMissions{} missions). Sources seldom peak at the same instant, but a
  budget built from per-source suprema must assume they do.
\end{enumerate}
The dependence gap can be closed only by knowing the coupling, which an
attacker who controls jamming denies the defender. The envelope gap shrinks
with better wind knowledge. The temporal gap, the largest of the three at
this $s$, can be closed without either, by bounding the supremum of the summed trajectory
directly, for example with a reachable-set computation per certification
window. Condition (C5), which licenses adding per-source suprema, held in
\nSubaddSources{} of joint missions.

\figGap

\section{Measured GNSS Marginal (RQ6)}
\label{sec:r5}
The pilot's GNSS marginal is simulated (99th percentile \nRFivePnnSim\,m). The
artifact contains measured navigation error from the real GPS-jamming flight of
the UAV Attack Dataset~\cite{whelan2020dataset}: \nRFiveSamples{} one-second
samples over a noise episode and a loss of fix, whose error grows with staleness
to \nRFiveEnvTen\,m at 10\,s and \nRFiveEnvMax\,m at 15\,s (spoofing samples
belong to $\dadv$ and are excluded). We replaced each simulated degradation
episode by the larger of its simulated deviation and the measured envelope at the
episode's staleness, and in a stress variant set every episode to the measured
45\,s duration. This never makes a mission's GNSS deviation smaller, but it
assumes, without measuring it, that navigation error becomes cross-track
deviation one for one.

At $m=10$\,m and 5\% FPR, with $\dadv=\nROneDadvFive$\,m, the floor estimated from
the empirical marginals falls from \nRFiveFloorSim{} (simulated GNSS) to
\nRFiveFloorMeas{} (measured envelope) and \nRFiveFloorLong{} (45\,s episodes).
With the finite-sample band of \cref{sec:revised}, the values are
\nRFiveFloorCpSim, \nRFiveFloorCpMeas{} and \nRFiveFloorCpLong. The bound is only
as good as its GNSS marginal; a simulated marginal overstated the certified
floor by a factor of two to three at this margin. It is the jamming-induced
tail, rather than the attack, that the corridor must absorb. At $m=20$\,m every variant
certifies \nRFiveFloorCpTwenty, the ceiling that \nCalMissions{} calibration
samples per source permit.

% ------------------------------------------------------------- moved from the body for length
\section{Negative Controls}
\label{app:negative}
A soundness test that cannot fail is uninformative, so we broke two conditions
deliberately and re-optimised the attacker under each (\cref{tab:negative}).
\begin{itemize}
\item \emph{(C2$'$) broken, no re-anchoring.} Undetected deviation reaches
  \nNegCtwoMaxDev\,m, and the test flags \nNegCtwoViol{} of \nNegCtwoRows{}
  points (\nNegCtwoViolFam{}), every one where the bound is informative.
\item \emph{(C4) broken, severe turbulence (45--60\,kt).} The fault mimic
  reaches \nNegCfourMaxDev\,m undetected, and the test flags the bound at
  $m=20$\,m, where $\pbar=0$ because the empirical marginals have finite
  support. At $m=10$\,m the observed frequency stays below $\pbar$, so this
  break is detected only where the bound is tightest. The weather mimic was
  detected in every mission, leaving no undetected missions to test.
\end{itemize}
The test thus detects unsoundness when a stated condition fails and enough
undetected missions remain to measure it, which is the behaviour the audit
loop of \cref{fig:framework} relies on.

\begin{table}[t]
\centering\footnotesize
\setlength{\tabcolsep}{2.4pt}
\caption{Negative controls at 5\% FPR. Undet.\ = undetected missions out of
\nEvalMissions; fails = failures among them; viol.\ = Holm-significant violation
of $\pbar$; dev.\ = largest undetected deviation (m).}
\label{tab:negative}
\begin{tabular}{@{}llcccccc@{}}
\toprule
\textbf{Control} & \textbf{disguise} & $m$ & \textbf{undet.} & \textbf{fails} & $\pbar$ & \textbf{viol.} & \textbf{dev.} \\
\midrule
\tabinput{tab_negative_body.tex}
\bottomrule
\end{tabular}
\end{table}

\section{Diagnosis of the Scaled-Test Violations (RQ5)}
\label{app:diagnosis}
With $\dadv=\nROneDadvFive$\,m, the test flagged \nROneViol{} of \nROneTests{}
points, all at $m=\nROneViolM$\,m, where the bound claimed $\pbar=0$. We
re-simulated each flagged cell with its seeds
(\code{r1\_violation\_diagnosis.csv}). The violations come from
\nDiagMissions{} missions, and the attacker caused none of them. Each fails
under its benign inputs alone (benign-only deviation \nDiagBenignMin--\nDiagBenignMax\,m),
the attack adds at most \nDiagAttackMax\,m, and \nDiagAllSat{} three saturate the
actuator with or without the attack. Two conditions of \cref{thm:floor} failed,
neither of which concerns disguise.
\begin{itemize}
\item \emph{(C3) for the marginals.} The empirical marginal of each source has
  finite support. Beyond the largest of \nCalMissions{} calibration samples,
  $\pbar$ is zero, and the bound claims a certainty that no finite sample can
  justify. One benign mission in 5{,}000 fails at $m=20$\,m.
\item \emph{(C1), which was not enforced.} Superposition held in at least
  \nROneSuperposMin{} of missions, and every violating mission saturated. The
  pilot had stated (C1) as an assumption, but no component of the stack checked
  it at run time.
\end{itemize}
This is a negative result about the pilot instantiation rather than about
\cref{thm:floor}. The theorem's conclusion requires its conditions, and the
scaled test found two that the instantiation did not secure. We report the
violations as found. The revision below fixes the instantiation before any new
data are examined and does not tune the old run until it passes.



\section{Detector Operating Curves}
\label{app:roc}
\Cref{fig:roc} shows the operating curves of the four detectors against a naive
5\,m step and a 0.5\,m/s ramp, which \cref{sec:eval} summarises.

\figROC

\section{PX4: the Pre-Registered Transfer Conditions}
\label{app:px4read}
The protocol fixed two conditions for a qualitative transfer of the disguise
matrix. (a)~With the gate armed, every primary attacker should evade in at least
half of the missions with at most twice its reduced-order deviation. After
deviations 4 and 5, \nPxReadAOk{} of five families met it, the step with
\nPxReadStepDev\,m against a limit of \nPxReadStepLim\,m. As registered,
\nPxReadAOkReg{} met it and the step reached \nPxReadStepDevReg\,m, a value that
came from a pair whose benign twin never flew. (b)~With the gate removed, EKF2's
gate should cap every family at \nPxReadCap\,m while D4 does not; D4 admitted
\nPxReadFaultDFour\,m to the fault mimic, but EKF2's gate also admitted
\nPxReadFaultDTwo\,m. What transferred is the ordering: the anchor gate decides
how far an undetected attacker gets, the inertial invariant D4 is bypassed by the
fault mimic, and without the gate the damage is two to three times larger in PX4
than in the reduced-order model.

\section{PX4: Deviations from the Protocol}
\label{app:px4dev}
All five are dated in the artifact's pre-registration, each recorded before the
data it changes were analysed, except deviation~5, which was recorded after the
stock-EKF2 results had been seen.
\begin{enumerate}
\item \emph{Validity rule.} Missions in which PX4 entered Descend or Termination
  before the window were excluded, with their pairs, and the thresholds were
  recalibrated on the \nPxBenign{} valid benign missions.
\item \emph{Re-selection of the PX4-optimised attackers.} The search had been
  steered by missions that never flew; the best attacker per family was
  re-selected on valid missions and evaluated on the pre-registered fresh seeds.
\item \emph{Robust statistics} (median and 95th percentile) next to every
  maximum.
\item \emph{Reruns.} All \nPxRerun{} missions excluded under deviation~1 were flown
  again with the same seeds, together with their pairs, and every one flew. Filling
  them in left the stealth budget at \nPxDadvReg\,m and the pre-registered reading
  unchanged.
\item \emph{Cruise condition.} The stock-EKF2 prediction at 5\,s failed on an
  attacked flight that never cruised and never applied its offset. It had passed
  the rule because PX4 entered mission mode inside the window before its
  failsafe. A mission is now valid only if it also cruised (at least 2\,m/s along
  track) before the attack onset, so that a slowdown caused by the attack is
  kept. The rule applies to every run: \nPxStillExcl{} missions of the v1.14
  study (\nPxStillExclPct) and \nStExcl{} of the replication are excluded. The
  stealth budget falls from \nPxDadvReg{} to \nPxDadv\,m. Every PX4 number in the
  body uses this rule. At the same time, the one-sided mismatch term of the body
  was added next to the registered two-sided one. That term,
  $\sup_{t\ge t_0}|e_{\mathrm{px4}}(t)-e_{\mathrm{lin}}(t)|$, also counts the linear
  loop overshooting PX4, and gives \nMmFivePointStat{} and \nMmOnePointStat\,m in
  place of \nMmFivePlusStat{} and \nMmOnePlusStat\,m.
\end{enumerate}
\begin{table}[t]
\centering\footnotesize
\caption{PX4 SITL, 5\,s re-anchoring, 5\% FPR, after deviations 4 and 5. Deviations
(m) are attack-only on undetected missions. RO: reduced-order study (detection: D2).}
\label{tab:px4}
\setlength{\tabcolsep}{1.5pt}
\begin{tabular}{@{}lrr@{}}
\toprule
 & PX4 & RO \\
\midrule
Step 5\,m caught: EKF2 / stack$+$gate & \nPxEOneStepDTwo{} / \nPxEOneStepAll & \nTprStepDTwo \\
Ramp 0.5\,m/s caught: EKF2 / stack$+$gate & \nPxEOneRampDTwo{} / \nPxEOneRampAll & \nTprRampDTwo \\
Transferred attackers: evasion & \nPxEvMin--\nPxEvMax & \nRoPrimEvMin--\nRoPrimEvMax \\
Transferred attackers: largest deviation & \nPxPrimMaxMin--\nPxPrimMaxMax & \nRoPrimMaxMin--\nRoPrimMaxMax \\
Worst of 35, stack$+$gate & \nPxWorstAllGate & \nAnchorMaxFive \\
Worst of 35, no gate: stack / D4 / D1 & \nPxWorstAllNoGate/\nPxWorstDFourNoGate/\nPxWorstDOneNoGate & \\
PX4-optimised: evasion; largest dev. & \nPxCemEvMin--\nPxCemEvMax; \nPxCemDevMax & \\
Stealth budget $\dadv$ & \nPxDadv & \nPxDadvRoPilot/\nPxDadvRoPooled \\
No anchor, ramp: evasion; drift & \nPxNoAnchorEv; \nPxNoAnchorDev & \\
Re-anchoring 1\,s: worst, stack$+$gate / D1$+$gate & \nPxTaWorstAllGate/\nPxTaWorstDOneGate & bound \nBoundTaOne \\
Stock EKF2: worst D1$+$gate, 5\,s / 1\,s & \nStDOneGateFive/\nStDOneGateOne & \nBoundTaFive/\nBoundTaOne \\
\bottomrule
\end{tabular}
\end{table}

\emph{Cause of the start-up failure.} The kept logs of all \nDiagFailed{} failed
attempts show invalid position setpoints right after take-off and a blind-land
failsafe. In \nDiagExec{} of them, the navigator began the mission without its
take-off step, which flights that succeeded rarely did. This points to a race
between the driver's switch to mission mode and PX4's take-off logic, independent
of the attack. It also occurs on v1.17.0, where \nStRetried{} of \nStMissions{}
missions needed a retry.

